% Sections 5.1 and 5.2, from lectures/L05/appendix/a_theory.md: one section per numbered section
% of the appendix.

\section{Dense Layer Architecture}
\label{c5:sec:architecture}\label{c5:app:a}

\subsection{Introduction}
\begin{itemize}
  \item A \textbf{dense layer} \emph{(fully connected layer)} is a layer in a neural network where
    every node receives all inputs from the previous layer.
  \item The class \code{ml::dense_layer::Dense} (see \secref{c5:sec:exercises}) implements the
    interface \code{ml::dense_layer::Interface} and represents such a layer, whether it's used as a
    hidden layer or output layer in the network.
\end{itemize}

\subsection{Structure of a layer}
A dense layer with \code{nodeCount} nodes and \code{weightCount} weights per node consists of the
parts described below. These map directly onto the theory in Chapter~\ref{c3:ch:nn} (see
\secref{c3:sec:nn}):

\begin{booktable}{@{}p{11em}p{12.4em}L@{}}
  \thead{Member variable} & \thead{Size} & \thead{Corresponds to (see \secref{c3:sec:nn})} \\
  \midrule
  \code{myOutput} & \code{nodeCount} & $y$, each node's output (after the activation function) \\
  \code{myPreActivationOutput} & \code{nodeCount} & $s$, each node's weighted sum (before the
    activation function) \\
  \code{myError} & \code{nodeCount} & $\Delta e$, each node's computed error \\
  \code{myBias} & \code{nodeCount} & $b$, each node's bias \\
  \code{myWeights} & \code{nodeCount} $\times$ \code{weightCount} & $w$, each node's weights \\
  \code{myActFunc} & none & $\sigma$, the layer's activation function \\
\end{booktable}

\subsection{Activation function}
Each layer is assigned an activation function of type \code{ActFunc} (\code{Relu}, \code{Tanh}, or
\code{None}) at construction. The activation function is applied to every node in the layer during
feedforward (implemented in the exercises).

\subsection{Composing a neural network}
A simple neural network (\code{ml::neural_network::Shallow}, see Chapter~\ref{c4:ch:network})
consists of two dense layers: a hidden layer and an output layer. The hidden layer's outputs are
the output layer's inputs, so the output layer's \code{weightCount} equals the hidden layer's
\code{nodeCount}.

% ------------------------------------------------------------------------------------------
\section{From Math to Code}
\label{c5:sec:math}
This section summarizes how the equations for feedforward, backpropagation, and optimization from
Chapter~\ref{c3:ch:nn} (see \secref{c3:sec:nn}) map onto the code you write in \code{Dense} (see
\secref{c5:sec:exercises} for the full instructions).

\subsection{Feedforward}
\[
  s = b + \sum_{i=0}^{n-1} w_i \cdot x_i, \qquad y = \sigma(s)
\]

\begin{narrowtable}{@{}ll@{}}
  \thead{Math} & \thead{Code} \\
  \midrule
  $b$ & \code{myBias[i]} \\
  $w_i$ & \code{myWeights[i][j]} \\
  $x_i$ & \code{input[j]} \\
  $s$ & \code{sum} (also stored in \code{myPreActivationOutput[i]}, see below) \\
  $\sigma(s)$ & \code{actFuncOutput(myActFunc, sum)} \\
  $y$ & \code{myOutput[i]} \\
\end{narrowtable}

\code{myPreActivationOutput[i]} is needed for backpropagation below, since the activation
function's derivative $y_p'$ must be computed from $s$ (the weighted sum before the activation
function was applied), not from $y$ (the output after the activation function).

\subsection{Backpropagation}
\textbf{Output layer:}
\[
  \delta = \yref - \yp, \qquad \Delta e = \delta \cdot y_p'
\]

\begin{narrowtable}{@{}ll@{}}
  \thead{Math} & \thead{Code} \\
  \midrule
  $\yref$ & \code{reference[i]} \\
  $\yp$ & \code{myOutput[i]} \\
  $\delta$ & \code{err} \\
  $y_p'$ & \code{actFuncDelta(myActFunc, myPreActivationOutput[i])} \\
  $\Delta e$ & \code{myError[i]} \\
\end{narrowtable}

\noindent\textbf{Hidden layer:}
\[
  \delta = \sum_{i=0}^{n-1} \left[\Delta e_i \cdot w_i\right], \qquad \Delta e = \delta \cdot y_p'
\]

\begin{narrowtable}{@{}ll@{}}
  \thead{Math} & \thead{Code} \\
  \midrule
  $\Delta e_i$ (next layer) & \code{nextLayer.error()[j]} \\
  $w_i$ (next layer) & \code{nextLayer.weights()[j][i]} \\
  $\delta$ & \code{err} \\
  $y_p'$ & \code{actFuncDelta(myActFunc, myPreActivationOutput[i])} \\
  $\Delta e$ & \code{myError[i]} \\
\end{narrowtable}

$\Delta e$ is the \emph{negative} gradient of the cost with respect to the node's weighted sum $s$,
which is why \code{optimize()} adds it rather than subtracting it. Most literature calls it the
\emph{delta} or the \emph{error term}. It's kept as \code{error()} here because that's what the
layer stores: one value per node.

It is not the same thing as the \code{inputGradients()} you'll meet in Chapter~\ref{c8:ch:conv}
(see the flatten layer in Chapter~\ref{c10:ch:flatten} for the clearest illustration of it), which
holds one value per \emph{input} and is what a layer hands back to the layer before it. A hidden
layer here fetches $\Delta e$ from its neighbour itself with \code{nextLayer.error()}, so no layer
ever needs to publish that second vector.

\subsection{Optimization}
\[
  \Delta c_n = \Delta e_n \cdot L, \qquad b_n = b_n + \Delta c_n, \qquad
  w_j = w_j + \Delta c_n \cdot y_j
\]

\begin{narrowtable}{@{}ll@{}}
  \thead{Math} & \thead{Code} \\
  \midrule
  $\Delta e_n$ & \code{myError[i]} \\
  $L$ & \code{learningRate} \\
  $b_n$ & \code{myBias[i]} \\
  $w_j$ & \code{myWeights[i][j]} \\
  $y_j$ (previous layer's output) & \code{input[j]} \\
\end{narrowtable}

\subsection{Helper functions \code{initRandGen()} and \code{randomStartVal()}}
These don't correspond to any specific equation, but handle the random initialization of weights
and biases described in Chapter~\ref{c3:ch:nn}. Both live in \file{ml/utils}, alongside the
activation-function helpers, rather than inside the layer itself:
\begin{itemize}
  \item \code{initRandGen()} ensures the random number generator (\code{std::rand()}) is seeded
    exactly once, regardless of how many layers are created. It's the same function you wrote in
    Chapter~\ref{c4:ch:network}, which is exactly why it lives in a shared file.
  \item \code{randomStartVal()} generates the actual random number in $[0.0, 1.0]$ that each bias
    and weight value is initialized with. Note this is a deliberately simple choice for the course:
    every parameter starts positive, with a mean of 0.5. Production networks initialize
    symmetrically around zero and scale by the number of inputs (Xavier or He initialization),
    which matters most for \code{Tanh} layers, since an all-positive start puts them straight into
    saturation where the gradients are small.
\end{itemize}
